% ============================================================================
% 中文注释：摘要（Abstract）—— 2026-11-03 按"贴膜玻璃是独立物理体制"重写
% ----------------------------------------------------------------------------
% 叙事顺序：
%   ① 普通 SIRR 的隐含假设：反射鲜明、透射清晰 ⇒ 主要矛盾是"分离"
%   ② 贴膜玻璃破坏了这个前提：膜是暗色的，两层同时被衰减 ⇒ 主要矛盾变成
%      "暗反射 + 低信噪透射 + 二者同频段" ⇒ 不是分离问题，而是辨识问题
%   ③ 由此得到四条可测失败模式（每条带数字）
%   ④ 采集层面的连带问题：黑布法引入系统性光照差 ⇒ 光度对齐与去反射必须统一
%   ⑤ 我们给的两件事：一个无掩码正则（+0.50 dB）+ 一套诚实的负结果与探针
% 关键数字：
%   反射/透射 = 0.201（Ω 外） / 0.112（Ω 内）；反射能量 34.5% 在高频
%   暗处反射扰动比亮处高 64%（23.06 vs 14.07 灰阶）
%   Ω 内 identity 24.52 dB 最高；公开方法 20.9–23.6；我们 22.39
%   Ω 内结构保有率：模型 0.48–0.65（GT=1）
%   逐图仿射 oracle 0.15–0.26 dB vs 留一场景共享 0.07 dB
%   L_quiet: +0.50 dB，95% CI [+0.34,+0.67]，14/15 场景为正
% 注意：Ω 内"透射被压暗"在 1019 面板上是 0.987（≈不变），只有 165530 是 0.830，
%       ⇒ 不要写成普适规律，写成"分场景、非系统性"。
% ============================================================================
\begin{abstract}
% 中文：普通 SIRR 假设"反射鲜明、透射清楚"，于是问题是分离。
Single-image reflection removal (SIRR) implicitly assumes a friendly photometric
regime: the reflected scene is bright, the transmitted scene is well exposed, and the
task is to \emph{separate} two comparable layers.

% 中文：贴膜玻璃破坏了这个前提——膜是暗色的，两层同时被压低。
We show that \textbf{tinted window film destroys this premise}. The film is dark: it
attenuates the reflected layer \emph{and} the transmitted layer at the same time. On a
hand-annotated $15$-scene panel shot through automotive window film we find that the film
treats the two layers \emph{differently}: the transmitted layer passes through the film bulk and
is scattered and dyed, while the reflection comes off the film surface. The result is a layer
asymmetry---the transmission is $2.3\times$ softer than on ordinary glass and the reflection is
$2.0\times$ sharper relative to it---which is \emph{not} explained by brightness: mean levels,
reflection magnitude and spectral overlap all match Real20. Under such a regime SIRR is not a
separation problem but an \emph{identification} problem.

% 中文：四条可测失败模式。
This reframing yields four measurable failure modes. (i)~The reflection is \emph{not}
strongest inside the occluder region $\Om$ but in the annulus around it: $4.3\times$ stronger than
inside $\Om$, covering $44\%$ of the frame, and that annulus is where every model is worst
($19.4$--$19.8\dBA$ vs.\ $23.5$ far away). Inside $\Om$ almost nothing needs to change yet every
method is worse than doing nothing (identity $24.52\dBA$ vs.\ $20.9$--$23.6$ for six public
methods and $22.4$ for ours), so the residual is read as a \emph{boundary}. (ii)~The damage is not a
photometric offset: after a per-image affine correction the models retain only
$0.48$--$0.65$ of the ground-truth structure inside $\Om$. (iii)~The reflection and the
scene share the same frequency band, so simply restoring the input's high frequencies
\emph{lowers} PSNR---faithful recovery requires semantic priors, not better filtering.
(iv)~What survives fine-tuning is per-image: a per-image affine oracle recovers
$0.15$--$0.26\dBA$ whereas a single shared affine fitted on other scenes recovers
$0.07\dBA$.

% 中文：采集层面的连带问题。
The standard ground-truth protocol---covering the glass with dark cloth---adds a
systematic illumination change, so photometric alignment and reflection removal cannot be
treated separately.

% 中文：我们给的东西 + 诚实边界。
We give a region-relative loss that is load-bearing ($-4.2\dBA$ if removed), a mask-free
change-aware regularizer derived from a random-disk control ($+0.50\dBA$, per-scene
$95\%$ CI $[+0.34,+0.67]$, positive on $14/15$ scenes, zero inference cost), and a
systematic set of negative results including perceptual losses, high-frequency splicing
and data mixing.
\end{abstract}
